\scanpage{77}
% manuscript line 1
réunion des $\varphi_{S'/S}(\mathcal U)$, où $U$ parcourt les ouverts affines de $\unsure{\mathfrak X|\mathscr X}$.

% manuscript line 3
\uline{Corollaire 2} Supposons \deleted{\uline{de plus}} \inserted{que tte partie finie de $\mathfrak X$ soit contenue dans un ouvert affine, et} $\mathfrak X$ de type fini \inserted{sur $S'$,} et $S$ loc. noeth. Alors $\varphi_{S'/S}(\mathfrak X)$ est de type fini sur $S$.

% manuscript line 6
La question étant locale \inserted{sur $S$}, on peut supposer $S$ affine. Tout résulte du

% manuscript line 8
\uline{Lemme} Soit $S$ noethérien, $\mathfrak X$ de type fini sur $S$, $n$ un entier. \inserted{\deleted{\unsure{les|Soit}} $(\mathfrak X_i)_{i\in I}$ famille d'ouverts,} On suppose que tte partie \deleted{finie} $E$ de $\mathfrak X$ telle que $\operatorname{card}(E)\leq n$ est contenue dans un ouvert \deleted{affine \unsure{de}} $\mathfrak X_i$. Alors il existe une \unsure{ss-famille|sous-famille} finie de $(\mathfrak X_i)_{i\in I}$ ayant la $=$ propriété.

% manuscript line 15
En effet, l'hyp. signifie que
% manuscript line 16
\[
\bigcup_{i\in I}
\underbrace{\mathfrak X_i\times_S\text{---}\times_S\mathfrak X_i}_{n\text{ facteurs}}
=
\underbrace{\mathfrak X\times_S\text{---}\times_S\mathfrak X}_{\unsure{n|\text{---}}\text{ facteurs}}.
\]
% manuscript line 17
On applique le fait que $\mathfrak X\times_S\text{---}\times_S\mathfrak X$ est \inserted{un espace} noethérien.

% manuscript line 19
$\unsure{\mathbb Z|2}$. Le foncteur $\underline W^{\unsure{\underline a|\wedge|a}}$ n'est pas local en tous les \unsure{cas}. Cas où il n'existe pas \unsure{même}
% archive sheet number: 76
